\section{Problemas de Valor de Contorno (PVC)}
Se da una ecuación diferencial, con las variables definidas en un dominio. Este representa
en sus bordes la frontera. De acá, se hacen $n+1$ nodos que dividen $n$ intervalos, los cuales
tienen todos la misma longitud, están equiespaciados. Esta longitud es \textbf{h}
\subsection{Fórmula genérica}
    $$\text{nodo } N-1: \quad \frac{U_N - 2U_{N-1} + U_{N-2}}{h^2} + U_{N-1} = 0$$
\subsection{Matriz tridiagonal}
Al aplicar diferencias finitas a la ecuación diferencial, obtenemos un sistema de ecuaciones evaluado en cada nodo interior:

\begin{align*}
    \text{nodo 1:} \quad & \frac{U_2 - 2U_1 + U_0}{h^2} + U_1 = 0 \\
    \text{nodo 2:} \quad & \frac{U_3 - 2U_2 + U_1}{h^2} + U_2 = 0 \\
    \vdots \quad & \\
    \text{nodo } N-1: \quad & \frac{U_N - 2U_{N-1} + U_{N-2}}{h^2} + U_{N-1} = 0
\end{align*}

\textit{Nota: $U_0$ y $U_N$ son datos (condiciones de contorno).}

\subsection{Sistema Matricial}
Agrupando términos, el sistema se puede expresar en forma matricial $A \cdot U = B$, resultando en una matriz tridiagonal:

\begin{equation}
\begin{bmatrix}
    -\frac{2}{h^2} + 1 & \frac{1}{h^2} & 0 & \dots & 0 \\[1em]
    \frac{1}{h^2} & -\frac{2}{h^2} + 1 & \frac{1}{h^2} & 0 \dots & 0 \\[1em]
    0 & \frac{1}{h^2} & -\frac{2}{h^2} + 1 & \frac{1}{h^2} \dots & 0 \\[1em]
    0 & 0 & \ddots & \ddots & \ddots
\end{bmatrix}
\begin{bmatrix}
    U_1 \\
    U_2 \\
    \vdots \\
    \vdots \\
    U_{N-1}
\end{bmatrix}
=
\begin{bmatrix}
    -U_0 / h^2 \\
    0 \\
    0 \\
    0 \\
    -U_N / h^2
\end{bmatrix}
\end{equation}

Como la matriz de coeficientes es tridiagonal, el sistema tiene altas chances de poder resolverse mediante métodos iterativos como Jacobi o Gauss-Seidel.